%------------------------------------------------------------------------------%
\section{Revisão}

\begin{exercises*}

% Exercício de prova do prof. Éden Amorim
\begin{exercise}
  Para cada função abaixo, apresente sua série de Taylor com o centro indicado.
  Apresente também o raio de convergência das séries.
  \begin{mathlist}
    \item q(x) = \frac{1}{5-x}      \), centrada em 4
    \item E(x) = \int e^{-x^2}\,dx  \), centrada em 0
    \item d(x) = \cos^2 x - \sin^2 x\), centrada em 0
  \end{mathlist}
  \begin{answer}
    \begin{alphalist}
      \item Pela série geométrica (razão $= x-4$):
      \[
      q(x) = \frac{1}{5-x}
      = \frac{1}{1-(x-4)}
      = \sum (x-4)^n
      \]
      para \(|x-4| < 1\)
      \item Pela série de Taylor da exponencial, substituindo $x \to -x^2$ e
      integrando a série (não afetam o raio, por ser infinito):
      \begin{align*}
        E(x) & = \int \left( \sum \frac{(-1)^n}{n!} x^{2n} \right) \, dx \\
        & = C + \sum \frac{(-1)^n}{(2n+1)n!} x^{2n+1}
      \end{align*}
      para $x \in \R$.
      \item Pela série de Taylor do cosseno, substituindo $x \to 2x$ (não afeta o raio, por ser infinito):
      \[
      d(x) =\cos(2x) = \sum \frac{(-1)^n 2^{2n}}{(2n)!}x^{2n}
      \]
      para $x \in \R$.
    \end{alphalist}
  \end{answer}
\end{exercise}

% Exercício de prova do prof. Éden Amorim
\begin{exercise}
  Considere a função
  \begin{align*}
    L(x) & = \ln \bigg( \frac{1+x}{1-x}\bigg) \\
    & = \ln(1+x) - \ln(1-x)
  \end{align*}
  para $\abs{x} < 1$
  Ela é utilizada para obter séries numéricas que aproximam os valores de $\ln 2$ e $\ln 3$.
  Nessa questão vamos explorar essa utilização.
  \begin{alphalist}
    \item Obtenha os valores de $x$ para os quais $L(x)=\ln 2$ e $L(x)=\ln 3$.
    \item Obtenha uma série de potências para $L$ centrada em zero, explicitando seu raio de convergência.
    \item Com base nos itens anteriores, apresente séries numéricas que convergem para $\ln 2$ e $\ln 3$
  \end{alphalist}
  \begin{answer}
    \begin{alphalist}
      \item Por cálculo direto, \(x=\sfrac{1}{3}\) e \(x=\sfrac{1}{2}\), respectivamente.
      \item Pela série de potências do logaritmo, substituindo $x \to - x$
      e somando as séries (não afetam o raio):
      \begin{align*}
        L(x) & = \sum \frac{1}{n}x^n - \sum \frac{(-1)^n}{n}x^n \\
        & = 2\sum \frac{1}{(2n+1)}x^{2n+1}
      \end{align*}
      para $|x| < 1$.
      \item Substituindo os valores obtidos no item (a) na série do item (b):
      \[\ln 2 = \sum \frac{2}{(2n+1)3^{2n+1}}\]
      \[\ln 3 = \sum \frac{1}{(2n+1)2^{2n}}\]
    \end{alphalist}
  \end{answer}
\end{exercise}

%------------------------------------------------------------------------------%

% Exercício de prova do prof. André Pereira
\begin{exercise} Seja $f(x)=\ln(x)$.
  \begin{alphalist}
    \item Mostre que \[f(x)=\sum_{n=1}^\infty\frac{(-1)^{n-1}(x-1)^n}{n}\] para \(x\in (0,2)\)
    \item Mostre que a série no item anterior converge absolutamente em $x\in (0,2)$.
    Analise se a série converge nos extremos $x=0$ e $x=2$.
    É possível afirmar que a série diverge em $(-\infty, 0)\cup (2, \infty)$?
    \item Some os quatro primeiros termos da série do item (a)
    para encontrar uma aproximação para o valor de $\ln(\sfrac{3}{2})$.
    Estime o erro cometido.
  \end{alphalist}
\end{exercise}

% Exercício de prova do prof. André Pereira
\begin{exercise}
  Seja \(f(x)=\frac{2}{1-x}\)
  \begin{alphalist}
    \item Encontre um padrão para $f^{(n)}(x)$ e use isso para construir
    a série de Taylor de $f$ centrada em $x=-1$.
    \item Mostre que
    \(f(x)=\frac{1}{1-\dfrac{x+1}{2}}\)
    \item
    Tomando \(r=\frac{x+1}{2}\), para quais valores de $x$
    a função $f$ pode ser vista como o limite da série geométrica \(\sum r^n\)?
    \item Use os itens (a) e (b) para concluir que $R_n(x)\rightarrow 0$
    se $x\in(-3, 1)$, onde $R_n(x)$ é o resto dado pelo teorema de Taylor.
  \end{alphalist}
\end{exercise}

% Exercício de prova do prof. André Pereira
\begin{exercise}
  Seja \(f(x)=e^{x}\)
  \begin{alphalist}
    \item Encontre a série de Taylor de $f$ centrada em $x=0$
    e mostre que $f$ coincide com a série em todo $x\in\R.$
    \item Escreva
    \[ \int_0^y f(-x^3)\,dx = \int_0^y e^{-x^3}\,dx \]
    como uma série de potências centrada em $y=0$.
    \item Calcule uma aproximação para
    \[ \int_0^{1/2} f(-x^3)\,dx \]
    fazendo a soma dos três primeiros termos da série obtida no item anterior.
    \item Estime o erro cometido na aproximação do item anterior.
  \end{alphalist}
\end{exercise}

% Exercício de prova do prof. André Pereira
\begin{exercise}
  Seja \( f(x)=\cos(x) \).
  \begin{alphalist}
    \item Encontre um padrão para $f^{(n)}(x)$ e use isso para encontrar a série de Taylor
    de $f$ centrada em $x=0$.
    Mostre que $f(x)$ coincide com sua série de Taylor para todo $x\in\R$.
    \item Escreva
    \[\int_0^x \cos(t^2)\,dt\]
    como uma série de potências centrada em $x=0$.
    \item Encontre uma aproximação para
    \[\int_0^1\cos(t^2)\,dt\]
    somando até o terceiro termo da série obtida no item anterior.
    Estime o erro cometido.
  \end{alphalist}
\end{exercise}

%    \begin{exercise}
%        \begin{mathlist}[2]
%            \item
%            \item
%            \item
%        \end{mathlist}
%        \begin{answer}
%            \begin{alphalist}[3]
%                \item
%                \item
%                \item
%            \end{alphalist}
%        \end{answer}
%    \end{exercise}

\end{exercises*}
%------------------------------------------------------------------------------%
